\documentclass{mynotes}

\begin{document}

\title{PGT theory derivation}
\author{Zhang Kai}
\date{\today}

\tableofcontents

\maketitle

\section{The analytical way of calculating PGT TFs}

\begin{equation}
	\begin{aligned}
		\widetilde{I}_{\text{i-s}}(\mathbf{k})
		= & P^*(\mathbf{k_{\text{inc}, \perp}})P(\mathbf{k}_{\perp}+\mathbf{k}_{\text{inc}, \perp})
		\frac{\pi \delta(k_z + k_{\text{inc}, z} - \sqrt{k_m^2 - \abs{\mathbf{k}_{\perp}+\mathbf{k}_{\text{inc}, \perp}}^2})}
		{\ii (k_z+k_{\text{inc},z})}
		\left[\widetilde{F}_{\text{re}}(\mathbf{k}) + \ii\widetilde{F}_{\text{im}}(\mathbf{k})\right]  \\
		  & -P(\mathbf{k_{\text{inc}, \perp}})P^*(\mathbf{k_{\text{inc}, \perp}}-\mathbf{k_{\perp}})
		\frac{\pi \delta(k_{\text{inc}, z} - k_z - \sqrt{k_m^2 - \abs{\mathbf{k}_{\text{inc}, \perp}-\mathbf{k}_{\perp}}^2})}
		{\ii (k_{\text{inc}, z} - k_z)}
		\left[\widetilde{F}_{\text{re}}(\mathbf{k}) - \ii\widetilde{F}_{\text{im}}(\mathbf{k})\right]. \\
	\end{aligned}
\end{equation}

Extract the first term

\begin{equation}
	\begin{aligned}
		\widetilde{I}_{\text{i-s}}(\mathbf{k})
		= & P^*(\mathbf{k_{\text{inc}, \perp}})P(\mathbf{k}_{\perp}+\mathbf{k}_{\text{inc}, \perp})
		\frac{\pi \delta(k_z + k_{\text{inc}, z} - \sqrt{k_m^2 - \abs{\mathbf{k}_{\perp}+\mathbf{k}_{\text{inc}, \perp}}^2})}
		{\ii (k_z+k_{\text{inc},z})}
	\end{aligned}
\end{equation}

The TFs can be assembled from

\begin{equation}
	\begin{aligned}
		\widetilde{I}_{\text{i-s}}(\mathbf{k})
		= & P^*(\mathbf{k_{\text{inc}, \perp}})P(\mathbf{k}_{\perp}+\mathbf{k}_{\text{inc}, \perp})
		\frac{\pi \delta(k_z + k_{\text{inc}, z} - \sqrt{k_m^2 - \abs{\mathbf{k}_{\perp}+\mathbf{k}_{\text{inc}, \perp}}^2})}
		{\ii (k_z+k_{\text{inc},z})}
	\end{aligned}
\end{equation}

TFs is accumulative, IFT to get the easy-to-assemble form

\begin{equation}
	\begin{aligned}
		P^*(\mathbf{k_{\text{inc}, \perp}})\ee^{-\ii \beta(k_{\text{inc}, \perp}) z}
		P(\mathbf{k}_{\perp}+\mathbf{k}_{\text{inc}, \perp})
		\frac{\ee^{\ii \beta(\mathbf{k}_{\perp}+\mathbf{k}_{\text{inc}, \perp}) z}}
		{2 \ii \beta(\mathbf{k}_{\perp}+\mathbf{k}_{\text{inc}, \perp})}
	\end{aligned}
\end{equation}
where $\beta(\mathbf{p}) = \sqrt{k_m^2 - \abs{\mathbf{p}}^2}$.

To assemble the final TFs, integrate the above equation over source

\begin{equation}
	\begin{aligned}
		\iint\,\dd^2 \mathbf{k}_{\text{inc}, \perp} S(\mathbf{k}_{\text{inc}, \perp})
		\left[ P(\mathbf{k_{\text{inc}, \perp}})\ee^{\ii \beta(k_{\text{inc}, \perp}) z} \right]^*
		\left[ P(\mathbf{k}_{\perp}+\mathbf{k}_{\text{inc}, \perp})
			\frac{\ee^{\ii \beta(\mathbf{k}_{\perp}+\mathbf{k}_{\text{inc}, \perp}) z}}
			{2 \ii \beta(\mathbf{k}_{\perp}+\mathbf{k}_{\text{inc}, \perp})} \right]
	\end{aligned}
\end{equation}

The above formula has the form of correlation and can be calculated
effectively by FFT

\begin{equation}
	\begin{aligned}
		A = \left[S(\mathbf{k}_{\text{inc}, \perp})
			P(\mathbf{k_{\text{inc}, \perp}})\ee^{\ii \beta(k_{\text{inc}, \perp}) z}\right]
		\otimes_{\text{corr}} \left[ P(\mathbf{k}_{\text{inc}, \perp})
			\frac{\ee^{\ii \beta(k_{\text{inc}, \perp}) z}}
			{2 \ii \beta(k_{\text{inc}, \perp})} \right]
	\end{aligned}
\end{equation}

Do the same thing with the second term

\begin{equation}
	\begin{aligned}
		B = \left[ P(\mathbf{k}_{\text{inc}, \perp})
			\frac{\ee^{\ii \beta(k_{\text{inc}, \perp}) z}}
			{2 \ii \beta(k_{\text{inc}, \perp})} \right]
		\otimes_{\text{corr}}
		\left[S(\mathbf{k}_{\text{inc}, \perp})
			P(\mathbf{k_{\text{inc}, \perp}})\ee^{\ii \beta(k_{\text{inc}, \perp}) z}\right]
	\end{aligned}
\end{equation}

FFT is the best way to calculate correlation

\begin{equation}
	\begin{aligned}
		F \otimes_{\text{corr}} G = \iint \dd\,\mathbf{x}' F^*(\mathbf{x}') G(\mathbf{x}+\mathbf{x}')
	\end{aligned}
\end{equation}

\begin{equation}
	\mathscr{F}\left[\iint \dd\,\mathbf{x}' F^*(\mathbf{x}') G(\mathbf{x}+\mathbf{x}')\right]
	=\mathscr{F}\left[F^*(-\mathbf{x}) \otimes_{\text{conv}} G(\mathbf{x})\right]
	=\widetilde{F}^*(\mathbf{k})\widetilde{G}(\mathbf{k})
\end{equation}

This equation says

\begin{equation}
	F(\mathbf{x}) \otimes_{\text{corr}} G(\mathbf{x})
	= \mathscr{F}^{-1}\left[\widetilde{F}^*(\mathbf{k})\widetilde{G}(\mathbf{k})\right]
	= \mathscr{F}^{-1}\left\{\mathscr{F}^*[F(\mathbf{x})]\mathscr{F}[G(\mathbf{x})]\right\}
\end{equation}

So, in order to calculate the second term, just take the conjugation
of the product of the two ffts, which explains how real/imaginary parts
correspond to the real/imaginary TFs.

The derivation above is casual and don't seem all that natural, but
at least it explains how Tian et al. calculated the TFs.
\end{document}
